% MCFG #2 v70. Standalone appendix section, included from main.tex.
% Adapted from the companion CFG paper, Appendix C.
\section{Substitutability of the bracketed doubling family}
\label{app:bracketed-substitutability}

This appendix supplies the self-contained proof of substitutability
used by Proposition~\ref{prop:fanout-one-lower-bound};
it follows~\cite[Appendix~C]{kuriyamaCFG}.
Write $L_i=L_{G_n}(E_i)$ for the grammar in that proposition. Then
$L_0=\{\mathtt a\}$ and
\[
 L_i=\lceil L_{i-1}\rceil \ \cup\
     \lfloor L_{i-1}L_{i-1}\rfloor
     \qquad (1\le i\le n).
\]
Every word of $L_i$ starts with precisely $i$ opening
brackets followed by $\mathtt a$.

A \emph{stack} is a word over
$\{E_0,\ldots,E_n,\rceil,\rfloor\}$ with top on the left.
Define partial transitions by
\[
\begin{aligned}
 E_0\psi&\xrightarrow{\mathtt a}\psi,&
 E_i\psi&\xrightarrow{\lceil}E_{i-1}\rceil\psi,\\
 E_i\psi&\xrightarrow{\lfloor}E_{i-1}E_{i-1}\rfloor\psi,&
 \rceil\psi&\xrightarrow{\rceil}\psi,\qquad
 \rfloor\psi\xrightarrow{\rfloor}\psi
 \qquad(1\le i\le n),
\end{aligned}
\]
and extend transitions to words. The system is deterministic;
each letter acts as an injective partial map on stacks.
Assign weights $\mathrm{wt}(E_i)=4^i$ and
$\mathrm{wt}(\rceil)=\mathrm{wt}(\rfloor)=1$.
Every transition strictly decreases the total weight.
A run on a nonempty top part $\alpha$ of $\alpha\psi$
therefore reaches $\psi$ at most once and never touches
the suffix $\psi$ before that point.

\begin{lemma}[Stack recognition of doubling blocks]
\label{lem:app-runs}
For $0\le i\le n$, every stack $\psi$ and $w\in\Gamma^*$,
\[
 E_i\psi\xrightarrow{w}\psi
 \quad\Longleftrightarrow\quad w\in L_i.
\]
If $\sigma\xrightarrow{w}\vartheta$ and $w\in L_i$,
then $\sigma=E_i\psi$ and $\vartheta=\psi$
for some stack $\psi$.
\end{lemma}

\begin{proof}
Induct on $i$. For $i=0$ the only possible return
reads $\mathtt a$. At $i\ge1$, the first letter
is $\lceil$ or $\lfloor$ and the resulting
exposed stack is respectively $E_{i-1}\rceil$
or $E_{i-1}E_{i-1}\rfloor$.
To return to $\psi$ one must exhaust the exposed
symbols in order; induction yields respectively
$\lceil z\rceil$ or $\lfloor z_1z_2\rfloor$
with $z,z_1,z_2\in L_{i-1}$. Conversely
every such word gives a returning run.
For the second assertion, the first $i$ opening
brackets of a word in $L_i$, followed by
$\mathtt a$, force the top symbols
$E_i,E_{i-1},\ldots,E_0$ in order;
the first assertion then identifies the
suffix left on the stack.
\end{proof}

\begin{lemma}[Branch synchronization]
\label{lem:app-sync}
If $\sigma\xrightarrow{x}\vartheta$ and
$\sigma\xrightarrow{y}\vartheta$ for nonempty
$x,y$ with different first letters, then
$\sigma=E_i\psi$ for some $i\ge1$, and
\[
 x=zx',\quad y=z'y',\quad z,z'\in L_i,
 \qquad
 \psi\xrightarrow{x'}\vartheta,\quad
 \psi\xrightarrow{y'}\vartheta.
\]
\end{lemma}

\begin{proof}
Only $E_i\psi$ with $i\ge1$ admits distinct
first letters, $\lceil$ and $\lfloor$.
Until the first visit to $\psi$, the
respective stacks have forms
$\alpha\rceil\psi$ and $\beta\rfloor\psi$;
their weights exceed $\mathrm{wt}(\psi)$,
while after such a visit weights are at most
$\mathrm{wt}(\psi)$.
Thus if just one of the two runs visited
$\psi$, their final weights would differ.
If neither visited $\psi$, their final
stacks would still contain incompatible
last exposed closing brackets.
Both runs must therefore visit $\psi$,
and Lemma~\ref{lem:app-runs} identifies the
prefixes read up to that visit.
\end{proof}

\begin{proposition}[Substitutability of the doubling family]
\label{prop:app-bracketed-substitutable}
For each $n\ge1$, if nonempty
$x,y\in\Gamma^+$ share a sentence context
in $L_n$, then
$\D_{L_n}(x)=\D_{L_n}(y)$.
\end{proposition}

\begin{proof}
First prove by induction on $|x|+|y|$
that if $\sigma\xrightarrow{x}\vartheta$
and $\sigma\xrightarrow{y}\vartheta$,
then $\D_{L_n}(x)=\D_{L_n}(y)$.
If one input is empty, strict weight decrease
forces both to be empty. If the first letters
agree, determinism gives the same next stack
and induction applies to the suffixes;
two-sided distributional equivalence is a
congruence, so the common prefix preserves it.
If they differ, Lemma~\ref{lem:app-sync}
writes $x=zx'$ and $y=z'y'$ with
$z,z'\in L_i$ and paths
$\psi\xrightarrow{x'}\vartheta$,
$\psi\xrightarrow{y'}\vartheta$.
Lemma~\ref{lem:app-runs} implies that
$z,z'$ induce identical partial stack
transformations, so
$\D_{L_n}(z)=\D_{L_n}(z')$.
Induction gives
$\D_{L_n}(x')=\D_{L_n}(y')$,
and congruence finishes the step.

Now if $uxv,uyv\in L_n$, determinism
gives a common stack $\sigma$ after $u$
in both accepting runs from $E_n$.
The stacks reached after $x$ and $y$
both read $v$ to the empty stack, hence
coincide by injectivity of the map induced by
$v$. The preceding claim proves the required
distributional equality.
\end{proof}
